%%
%% This is file `poli-max-exemplo.tex',
%% generated with the docstrip utility.
%%
%% The original source files were:
%%
%% ufrj-poli.dtx  (with options: `polimaxexemplo')
%% 
%% This is a sample undergraduate project which illustrates the use of the
%% `ufrj' document class with the `ufrj-poli' unit style. It is generated from
%% ufrj-poli.dtx by docstrip.
%% 
%% Copyright (C) 2008-2026 Vicente Helano F. Batista, George O. Ainsworth Jr.,
%% Geraldo Xexéo and Eduardo Mangeli. Released under the GNU General Public
%% License version 3 (see the COPYING.txt file).
%% 
\documentclass[grad]{ufrj}
\usepackage[times]{ufrj-poli}

\usepackage{graphicx}
\usepackage{makeidx}

\addbibresource{exemplo.bib}

\makelosymbols
\makeloabbreviations
\makelosiglas
\makeglossarylist
\makeindex

\defineacronym{ufrj}{UFRJ}{Universidade Federal do Rio de Janeiro}
\defineacronym{poli}{Poli}{Escola Politecnica}

\begin{document}
  \title{Titulo do Projeto de Graduacao}
  \foreigntitle{Title of the Undergraduate Project}
  \subtitle{subtitulo, quando houver}
  \foreignsubtitle{subtitle, if any}
  \author{Nome do Autor}{Sobrenome}
  \advisor{Nome do Orientador}{Sobrenome}{D.Sc.}{UFRJ}
  \coadvisor{Nome do Coorientador}{Sobrenome}{D.Sc.}{UFRJ}
  % A banca e o que se declara aqui, na ordem: o orientador, que a preside
  % (3.1.2.1.3e), e o primeiro.
  \examiner{Nome do Orientador Sobrenome}{D.Sc.}{UFRJ}
  \examiner{Nome do Coorientador Sobrenome}{D.Sc.}{UFRJ}
  \examiner{Nome do Primeiro Examinador Sobrenome}{D.Sc.}{UFRJ}
  \examiner{Nome do Segundo Examinador Sobrenome}{Ph.D.}{UFRJ}
  % A sigla do curso: ECI, CIVIL, AMBIENTAL, MECANICA, NAVAL, PRODUCAO,
  % ELETRICA, ELETRONICA, CONTROLE, MATERIAIS, NUCLEAR, PETROLEO, QUIMICA ou
  % NANOTECNOLOGIA. Com a opcao civilpordepartamento, as cinco siglas de
  % departamento da Engenharia Civil tambem valem.
  \department{ECI}
  \date{01}{2026}
  \approvaldate{15 de janeiro de 2026}
  % A cidade e da CLASSE: o estilo declara o Rio de Janeiro, e o autor troca
  % com \city (ou \place) se o trabalho for feito em Macae ou em Duque de
  % Caxias.
  \keyword{primeira palavra-chave}
  \keyword{segunda palavra-chave}
  \foreignkeyword{first keyword}
  \foreignkeyword{second keyword}

  \maketitle
  \frontmatter

  \dedication{A quem leu a norma antes de escrever.}

  \epigrafe{``A citacao que abre o trabalho, com a mesma
    formatacao da citacao longa.''}{Autor da epigrafe}

  \begin{abstract}
    O resumo em portugues, em paragrafo unico, com 150 a 500 palavras
    (3.1.2.1.4). Ele nao repete o titulo, nao cita referencia e nao traz
    formula: diz o que foi feito, como, e o que se concluiu. A folha traz a
    referencia do proprio trabalho no alto e as palavras-chave no pe, como
    pedem os Anexos E e F do Manual.
  \end{abstract}

  \begin{foreignabstract}
    The abstract in English, the version of the abstract in the language of
    international dissemination (3.1.2.1.5), with the same content as the
    Portuguese one.
  \end{foreignabstract}

  % As listas, na ordem da 3.1.2: ilustracoes (figuras, quadros), tabelas,
  % abreviaturas, siglas e simbolos. O sumario e sempre o ultimo, e a classe
  % confere a ordem.
  \listoffigures
  \listofquadros
  \listoftables
  \printloabbreviations
  \printlosiglas
  \printlosymbols
  \tableofcontents

  \mainmatter

  \chapter{Introducao}
  \label{cap:introducao}
  O texto do trabalho comeca aqui. A \sigla{poli} e uma unidade da
  \sigla{ufrj}, e na segunda vez a sigla sai sozinha: a \sigla{poli}.

  A citacao segue o sistema da classe \cite{m-4211}, e a citacao com o autor
  no texto sai como em \textcite{m-4212}. A citacao de mais de tres linhas vai
  no ambiente proprio:

  \begin{longquote}
    A citacao longa entra com recuo de 4 cm da margem esquerda, em corpo
    menor, sem aspas e com espacamento simples, como pede a 2.5 do Manual.
    Depois dela, o texto continua no paragrafo seguinte.
  \end{longquote}

  \section{Uma secao}
  O simbolo \symbl{$\alpha$}{coeficiente de dilatacao} entra na lista de
  simbolos na ordem em que aparece, e a abreviatura \abbrev{Prof.}{Professor}
  e o \acron{PG}{Projeto de Graduacao} entram nas suas. O
  \glossaryterm{termo tecnico}{a definicao do termo, que sai no glossario}
  entra no glossario, e esta palavra\index{palavra} entra no indice
  remissivo.

  \chapter{Ilustracoes, quadros e tabelas}
  \label{cap:ilustracoes}

  A legenda vem em cima e a fonte embaixo, as duas alinhadas a esquerda, na
  margem da ilustracao (2.10 e Figura 1 do Manual).

  \begin{figure}[htbp]
    \caption{Uma figura estreita, com a legenda na largura dela}
    \illustrationwidth{6cm}
    \centering
    \rule{6cm}{3cm}
    \source{Elaborado pelo autor (2026).}
  \end{figure}

  \begin{quadro}[htbp]
    \caption{Um quadro: informacao textual, com as quatro bordas}
    \centering
    \begin{tabular}{|l|l|}
      \hline
      Elemento & Onde esta \\
      \hline
      Capa & 3.1.1.1 \\
      Folha de rosto & 3.1.2.1.1 \\
      \hline
    \end{tabular}
    \source{Elaborado pelo autor (2026).}
  \end{quadro}

  \begin{table}[htbp]
    \caption{Uma tabela: dado numerico, sem as bordas laterais}
    \centering
    \begin{tabular}{lr}
      \hline
      Curso & Alunos \\
      \hline
      Engenharia de Computacao e Informacao & 100 \\
      Engenharia Civil & 200 \\
      \hline
    \end{tabular}
    \source{IBGE (2026).}
  \end{table}

  \chapter{Listagem de programa}
  \label{cap:programa}
  A listagem sai com o estilo da linguagem, e o flutuante dela tem lista
  propria:

  \begin{python}[caption={Um programa de exemplo},label={prog:exemplo}]
def soma(a, b):
    return a + b
  \end{python}

  \chapter{Conclusao}
  \label{cap:conclusao}
  A conclusao do trabalho, que retoma o que o Capitulo~\ref{cap:introducao}
  propos.

  \appendix
  \chapter{Um apendice}
  O apendice e texto do proprio autor, complementar ao trabalho.

  \annex
  \chapter{Um anexo}
  O anexo e documento de terceiro, trazido para o trabalho.

  \backmatter
  \printbibliography
  \printglossarylist
  \printindex
\end{document}
%% 
%%
%% End of file `poli-max-exemplo.tex'.
